\documentclass[12pt,a4paper]{article} 

\usepackage[utf8]{inputenc} 
\usepackage[polski]{babel} 
\usepackage[T1]{fontenc} 
\usepackage{geometry} 
\geometry{ 
    top=2.5cm, 
    bottom=2.5cm, 
    left=2.5cm, 
    right=2.5cm 
} 
\usepackage{enumitem} 
\usepackage{hyperref} 
\usepackage{tikz} 
\usetikzlibrary{shapes.geometric, arrows, positioning} 

\renewcommand{\contentsname}{Spis treści} 

\begin{document} 

% -- STRONA TYTUŁOWA (Strona 1) -- 
\begin{titlepage} 
    \begin{center} 
        \vspace*{1cm} 
        
        {\LARGE \textbf{Politechnika Warszawska}} \\[0.2cm]  
        {\Large Wydział Elektryczny} \\[2cm]  
        
        {\large Kierunek: Informatyka stosowana} \\[0.1cm]  
        {\large Przedmiot: Języki i metody programowania 2} \\[3.5cm]  
        
        {\Huge \textbf{Specyfikacja Implementacyjna}} \\[1cm]  
        {\Large Projekt C: Wyznaczanie współrzędnych węzłów dla wizualizacji\\ grafu planarnego} \\[3.5cm]  
        
        {\large \textbf{Autorzy:}} \\[0.1cm]  
        {\large Patryk Pawlak 343367, Krzysztof Dębski 343317}  
        
        \vfill 
        
        {\large Warszawa, 22 marca 2026}  
        \vspace*{0.5cm} 
    \end{center} 
\end{titlepage} 

\newpage 

% -- SPIS TREŚCI (Strona 2) -- 
\tableofcontents 
\newpage 

% -- TREŚĆ DOKUMENTU (Strona 3) -- 

\section{Architektura systemu}

Aplikacja została zaprojektowana w oparciu o architekturę modularną, dzielącą system na niezależne komponenty odpowiedzialne za konkretne zadania. Wyodrębniono następujące warstwy: 

\begin{itemize} 
    \item \textbf{Warstwa CLI:} Odpowiada za obsługę interfejsu linii poleceń i parsowanie flag wejściowych bezpośrednio w module głównym aplikacji. 
    \item \textbf{Warstwa I/O:} Zajmuje się wczytywaniem pliku tekstowego z listą krawędzi oraz zapisem wygenerowanych współrzędnych do pliku w wybranym formacie. 
    \item \textbf{Warstwa Logiki:} Główny silnik obliczeniowy zawierający implementację algorytmu Fruchterman-Reingold oraz metody bazującej na Twierdzeniu Tutte'a. 
\end{itemize} 

\vspace{1cm} 
\begin{center} 
    \begin{tikzpicture}[node distance=1.5cm, auto] 
        \tikzstyle{block} = [draw, rectangle, fill=orange!10, minimum height=3.5em, minimum width=9em, text width=8.5em, text centered] 
        
        \node [block] (main) {core/ (Główny sterownik CLI)}; 
        \node [block, below left=of main, xshift=-0.5cm, yshift=-1cm] (io) {io/ (Obsługa struktur i plików)}; 
        \node [block, below right=of main, xshift=0.5cm, yshift=-1cm] (algo) {algorithms/ (Logika rozkładu)}; 
        \node [block, below=of algo, yshift=-0.5cm] (utils) {utils/ (Funkcje mat.)}; 
        
        \path [draw, thick, ->, >=stealth] (main) -- (io); 
        \path [draw, thick, ->, >=stealth] (main) -- (algo); 
        \path [draw, thick, ->, >=stealth] (io) -- (algo); 
        \path [draw, thick, ->, >=stealth] (algo) -- (utils); 
    \end{tikzpicture} 
    \\ \vspace{0.4cm} 
    \textbf{Schemat 1:} Diagram architektury modułowej i przepływu danych. 
\end{center} 

\vspace{1cm} 

\section{Technologie}

Aplikacja wykorzystuje następujące technologie i narzędzia: 

\begin{itemize} 
    \item \textbf{Język programowania:} C. 
    \item \textbf{Kompilator:} GCC. Proces kompilacji jest zarządzany automatycznie z wykorzystaniem programu \texttt{make} oraz pliku \texttt{Makefile} z flagami \texttt{-Wall -pedantic}. 
    \item \textbf{Narzędzia:} Valgrind do monitorowania poprawności alokacji i wykrywania ewentualnych wycieków pamięci podczas pracy z grafem. 
\end{itemize} 

\newpage 
% -- STRUKTURA I INTERFEJSY (Strona 4) -- 

\section{Struktura projektu}

Zasoby projektu, w tym kod źródłowy oraz dokumentacja, zostały zorganizowane w strukturze katalogów zgodnej z repozytorium \texttt{src\_c}: 

\begin{itemize} 
    \item \texttt{doc/} -- katalog przeznaczony na pełną dokumentację projektu. 
    \item \texttt{src\_c/algorithms/} -- folder zawierający implementacje algorytmów rozmieszczania węzłów grafu na płaszczyźnie. 
    \item \texttt{src\_c/core/} -- moduł główny programu odpowiedzialny za logikę sterowania aplikacji oraz obsługę interfejsu użytkownika. 
    \item \texttt{src\_c/io/} -- katalog grupujący funkcjonalności związane z definicją struktur danych grafu oraz operacjami wczytywania i zapisywania plików. 
    \item \texttt{src\_c/utils/} -- zbiór pomocniczych bibliotek matematycznych oraz funkcji wspierających główne obliczenia algorytmiczne. 
\end{itemize} 

\vspace{1.5cm} 

\section{Interfejsy i Specyfikacja API}

System komunikuje się za pomocą precyzyjnie określonych interfejsów, tworzących spójne API. Zdefiniowano następujące parametry, odpowiedzi i kody operacji: 

\begin{itemize} 
    \item \textbf{Interfejs linii poleceń (CLI):} Program obsługuje parametry uruchomieniowe (polecenia): \texttt{-i} (ścieżka wejściowa), \texttt{-o} (ścieżka wyjściowa), \texttt{-a} (wybór algorytmu), \texttt{-n} (liczba iteracji), \texttt{-b} (format binarny), \texttt{-v} (tryb szczegółowy) oraz \texttt{-h} (pomoc). 
    \item \textbf{Formaty danych (I/O):}  
    \begin{itemize} 
        \item \textit{Parametry wejściowe z pliku:} \texttt{<nazwa\_krawędzi> <wierzchołek\_A> <wierzchołek\_B> <waga>}. 
        \item \textit{Odpowiedź wyjściowa w pliku:} \texttt{<wierzchołek> <x> <y>}. 
    \end{itemize} 
    \item \textbf{Specyfikacja wewnętrznego API (zdefiniowane w \texttt{src\_c/io/graph.h}):} 
    \begin{itemize} 
        \item \texttt{Graph *create\_graph(void)} -- polecenie inicjalizacji i alokacji struktur grafu. Zwraca wskaźnik na graf. 
        \item \texttt{int graph\_load\_text(Graph *g, const char *path, int verbose)} -- przyjmuje ścieżkę do pliku jako parametr i ładuje dane. Zwraca kody błędów w razie niepowodzenia. 
        \item \texttt{int graph\_save\_text(const Graph *g, const char *path)} -- polecenie eksportu danych do formatu tekstowego. 
        \item \texttt{int graph\_save\_binary(const Graph *g, const char *path)} -- polecenie eksportu danych do formatu binarnego. 
        \item \texttt{void free\_graph(Graph *g)} -- polecenie bezpiecznego zwalniania wszystkich zasobów pamięciowych powiązanych z parametrem \texttt{g}. 
    \end{itemize} 
\end{itemize} 

\newpage 
% -- MODEL DANYCH I ZASADY (Strona 5) -- 

\section{Struktura danych}

Model danych opiera się na powiązanych strukturach dynamicznych, które odwzorowują (jako encje i relacje) elementy grafu planarnego: 

\begin{itemize} 
    \item \texttt{Node} -- encja pojedynczego węzła (identyfikator oraz współrzędne położenia). 
    \item \texttt{Edge} -- encja krawędzi. Stanowi fizyczną relację łączącą węzły końcowe (zawiera indeksy do tablicy węzłów, nazwę i wagę). 
    \item \texttt{Graph} -- główny agregat zarządzający dynamicznymi tablicami encji \texttt{Node} i \texttt{Edge}. 
    \item \texttt{Config} -- struktura pomocnicza agregująca parametry działania programu przekazanych przez użytkownika. 
\end{itemize} 

\vspace{1cm} 
\begin{center} 
    \begin{tikzpicture}[node distance=2.5cm, auto] 
        \tikzstyle{entity} = [draw, rectangle, fill=green!10, minimum height=3.5em, minimum width=9em, text centered, font=\bfseries] 
        
        \node [entity] (graph) {Struktura Graph}; 
        \node [entity, right=of graph, xshift=1.5cm] (node) {Encja Node}; 
        \node [entity, below=of graph, yshift=-0.5cm] (edge) {Encja Edge}; 
        
        \path [draw, thick] (graph) -- node[above] {1..n} (node); 
        \path [draw, thick] (graph) -- node[left] {1..n} (edge); 
        \path [draw, thick] (edge) -| node[pos=0.7, right] {łączy 2} (node); 
    \end{tikzpicture} 
    \\ \vspace{0.4cm} 
    \textbf{Schemat 2:} Model relacji danych (ERD). 
\end{center} 

\vspace{1cm} 

\section{Zasady implementacji}

W celu zapewnienia stabilności i czytelności kodu przyjęto następujące standardy: 

\begin{itemize} 
    \item \textbf{Standardy kodowania:} Restrykcyjne przestrzeganie zasad DRY (unikanie duplikacji), KISS (prostota kodu) oraz YAGNI (nie implementujemy na zapas). Cały stan aplikacji przekazywany jest jawnie poprzez wskaźniki (Minimal Global State), co eliminuje potrzebę stosowania zmiennych globalnych. 
    \item \textbf{Zarządzanie pamięcią:} Projekt zakłada bezwzględną weryfikację poprawności każdej operacji alokacji pamięci. Program każdorazowo sprawdza, czy funkcje takie jak \texttt{malloc} nie zwróciły wartości \texttt{NULL}, a każda zaalokowana struktura posiada odpowiadający jej mechanizm zwalniania pamięci.
    \item \textbf{Programowanie defensywne i kody błędów:} Wszystkie operacje na plikach i danych wejściowych są walidowane. Błędy skutkują zakończeniem programu ze zdefiniowanym w API kodem: \texttt{ERR\_FILE} (1), \texttt{ERR\_FORMAT} (2), \texttt{ERR\_ALGO} (3) lub \texttt{ERR\_ARGS} (4). 
    \item \textbf{Testowanie:} Weryfikacja poprawności działania odbywa się poprzez scenariusze testowe nadzorowane przez narzędzie Valgrind. 
\end{itemize} 

\end{document}
